\chapter{Balazard, Saias and Yor meet Thorin}\label{chap:bsy}\label{bsy:sec:intro}\chaptermark{Balazard, Saias and Yor meet Thorin}

Balazard, Saias and Yor \cite{BSY99} proved in 1999 that the Riemann Hypothesis is equivalent to a single Cauchy average of $\log|\zeta|$ on the critical line being zero. The Thorin programme of this book studies the measures $\mathfrak U_\sigma$ with density $m_\sigma=\pi^{-1}\Re\,\xi'/\xi(\sigma+it)$ and their images $U_\sigma$, the Thorin measures of $\xi(\sigma)/\xi(\sigma+\sqrt s)$ (equation~\eqref{lo:eq:Utheta} and Chapter~\ref{chap:lo}). This chapter shows that the two are the same harmonic analysis, one derivative in $\sigma$ and one archimedean term apart, and then solves the resulting Cauchy averages numerically by quadrature, by plain Monte Carlo and by a vertical likelihood sampler after Polson and Scott (Section~\ref{bsy:sec:mc}; computed, where the plain sampler has the smaller standard error). Chapter~\ref{chap:cf} is the companion: it cites Theorem~\ref{bsy:thm:bsy} for the criterion, and its particle filter and nested sampling over heights supply the scans of Section~\ref{bsy:sec:mc}.

The chapter is based on \cite{P26bsy}. The zeros and counts behind Theorem~\ref{bsy:thm:dhid} are certified in ball arithmetic; the other tables were recomputed in floating point, and where the recomputation differed from the note the recomputed value is printed.

Every value at the level $\frac78$ is conditional on the quasi Riemann hypothesis $\mathrm{QRH}(\frac78)$, announced by OpenAI \cite{OAI26a}, formalised in Lean and not refereed (Section~\ref{lv:sec:qrh}). The notation is that of Chapters~\ref{chap:lo} to~\ref{chap:cf}: $\mathfrak U_\sigma$ is the measure on $\R$ without the factor $\frac12$; $P_a(x)=a/(\pi(a^2+x^2))$ is the Cauchy density of scale $a$ and, for $a<0$, the same formula, so that $P_{-a}=-P_a$ as in Proposition~\ref{prop:struct}; and $\mathrm{QRH}(\theta)$ is the statement that $\zeta$ has no zero with $\beta>\theta$.

The chapter contains:
\begin{enumerate}\setlength{\itemsep}{1pt}
\item a dictionary between the logarithmic level (Balazard, Saias and Yor) and the Thorin level: the Poisson integral becomes a Cauchy smoothing of $\mathfrak U_{1/2}$, and the Green potential of each zero off the line a signed Cauchy defect, with an archimedean term that the logarithmic level does not see (Section~\ref{bsy:sec:dictionary});
\item the criterion at every level $\theta\in[\frac12,1)$ and every height, as the Nevanlinna formula of the half plane for $\zeta(s)(s-1)/(s-2\theta+1)$, which characterises $\mathrm{QRH}(\theta)$; the case $\theta=\frac12$, $t_0=0$ is the theorem of Balazard, Saias and Yor (Section~\ref{bsy:sec:bsy});
\item the reading of the 2 $\sin^2$ lemma (equation~\eqref{eq:sin2}): the L\'evy density of the reciprocal Xi function is positive with no hypothesis, and RH is equivalent to that reciprocal being an extended generalised gamma convolution (EGGC) with the Thorin measure $\mathfrak U_{1/2}$ (Section~\ref{bsy:sec:sin2});
\item a probabilistic reading of the criterion by Brownian motion and the Riesz decomposition, in the spirit of Yor (Section~\ref{bsy:sec:bm}), and its complex form: under $\mathrm{QRH}(\theta)$ the boundary values of $\log h_\theta$ are the stopped value of a time changed complex Brownian motion, so $\log|\zeta|$ and $\arg h_\theta$ have equal variances and zero covariance under harmonic measure at every height; without the hypothesis the defects of these moments are integrals along horizontal cuts from the zeros to the line (Proposition~\ref{bsy:prop:cuts});
\item a proof that the 2 $\sin^2$ L\'evy density of the Davenport and Heilbronn function is positive although it has zeros off the line, the three finite zero computations behind it being certified in ball arithmetic (Section~\ref{bsy:sec:sin2});
\item a vertical (layer cake) form of the criterion, in which RH says that the upper and lower tails of $\log|\zeta|$ under harmonic measure balance, solved by quadrature, by plain and vertical Monte Carlo, and by the weighted Monte Carlo of Yakowitz, Krimmel and Szidarovszky with an exact Dirichlet polynomial control; an exact local form of the criterion on a half disc, for which the sorted trapezoid rule has mean square error $O(N^{-4})$, and $O(N^{-5})$ after a change of variable (proved); and the truncated whole line average, whose tail is $O(T^{-2}\log T)$ under RH \cite{BLM14} (computed in Table~\ref{bsy:tab:trunc}; Section~\ref{bsy:sec:mc});
\item a smoothed Riesz identity in the plane, with the Cauchy kernel replaced by a compactly supported bump. The integrand is then smooth away from the zeros, the Voronoi and Delaunay forms of the weighted rule have root mean square error $O(N^{-3/2})$ at independent uniform nodes (proved, Theorem~\ref{bsy:thm:wmc}, which gives $O(N^{-1/2-2/d})$ in dimension $d$; computed at jittered nodes), a shifted grid reaches machine precision in a box free of zeros, and deflation handles a zero in the box; the same count is a linear functional of the horizontal logarithmic derivative alone (Section~\ref{bsy:sec:2d}).
\end{enumerate}
Assuming RH, Bui, Lester and Milinovich \cite{BLM14} bound the tail of the truncated Balazard, Saias and Yor integral and disprove a conjecture of Borwein, Bradley and Crandall \cite{BBC00} on its rate of decay; Section~\ref{bsy:sec:trunc} uses their results. Except for the positivity of $\mathfrak U_\sigma$ for $\sigma\ge1$ and Theorem~\ref{bsy:thm:cmart}, every identity of Sections~\ref{bsy:sec:thorin} to~\ref{bsy:sec:bm} holds, with the changes listed in Section~\ref{bsy:sec:limits}, for the Davenport and Heilbronn function, which has zeros off the line; Section~\ref{bsy:sec:limits} records what that implies.
